%% Section: adequacy, input by paper.tex after proof.tex. Sources: agda/Mu/Theorems/Adequacy.agda (the
%% statements), agda/Mu/proof/Adequacy/{FPC,Soundness,LogRel,X86*,
%% EndToEnd}.agda. Size from loc.md (adequacy table, interpreters).

\section{Adequacy}
\label{sec:adequacy}

Theorem~\ref{thm:compiler} speaks of denotations. To
connect it to running code, we prove that the graph-model semantics of \fpc\ and \instr\ are adequate for the clocked interpreters of Section~\ref{sec:evaluation} (Appendix~\ref{app:interp}).

\begin{theorem}[Adequacy]\label{thm:adequacy}
Let $M$ be a closed term of \fpc. For every constant $c$,
$c \in \den{M}$ if and only if the interpreter of \fpc\ terminates on
$M$ with $c$, and $\den{M}$ is empty if and only if the interpreter
diverges on $M$, that is, times out at every clock. The same holds
for \instr\ and $\pass{compile}(M)$.
\end{theorem}

The proof for each language has two halves. Soundness says that if
the interpreter returns a value at some clock, then that value denotes
a part of the meaning of the program. Adequacy says that every element of the meaning is
related to the value that the interpreter returns from some clock on,
by a logical relation defined by recursion on the element, as in
earlier adequacy proofs for graph and filter models
\citep{Siek:2017ab,Wadler:2020aa}. Together with
Theorem~\ref{thm:compiler}, Theorem~\ref{thm:adequacy} gives the
theorem that the operational proofs state, namely that for a program
of integer type the interpreters of \fpc\ and \instr\ terminate with
the same integer or both diverge. The adequacy proofs are 1,529 lines of Agda; Section~\ref{sec:evaluation} explains why we report them separately.
